\documentclass[11pt,a4paper]{article}
\usepackage[margin=1in]{geometry}
\usepackage{amsmath,amssymb,amsfonts,amsthm}
\usepackage{bm}
\usepackage{booktabs}
\usepackage{enumitem}
\usepackage{hyperref}
\usepackage{cleveref}

% ── Notation macros (consistent with Pythonic-DISORT.ipynb) ──────────────
\newcommand{\R}{\mathbb{R}}
\newcommand{\diag}{\operatorname{diag}}
\newcommand{\vect}[1]{\bm{#1}}
\newcommand{\mat}[1]{\mathbf{#1}}
\newcommand{\tbot}{\tau_{\mathrm{bot}}}
\newcommand{\dd}{\,\mathrm{d}}
\newcommand{\tr}{\operatorname{tr}}
\newcommand{\cond}{\operatorname{cond}}
\newcommand{\avec}{\operatorname{vec}}

\theoremstyle{definition}
\newtheorem{proposition}{Proposition}
\newtheorem{remark}{Remark}
\newtheorem{definition}{Definition}

\title{Invariant-Imbedding Riccati Solver for Radiative Transfer\\
with Continuously $\tau$-Varying Optical Properties}
\author{Technical Report --- Extension to the Pythonic-DISORT\\
Comprehensive Documentation}
\date{March 2026}

\begin{document}
\maketitle

\begin{abstract}
This report documents the Riccati forward solver implemented in
\texttt{pydisort\_riccati\_jax}, an alternative entry point to PythonicDISORT
for atmospheres where the single-scattering albedo $\omega(\tau)$ and
phase function vary continuously with optical depth.  Instead of
per-layer eigendecomposition, the solver integrates the
invariant-imbedding Riccati ODE for reflection and transmission
operators using diffrax's Kvaerno5 method~\cite{kidger2022neural}
(5-stage ESDIRK, L-stable, order~5) with PIDController adaptive
step-size control, running on JAX~\cite{jax2018github} for automatic
differentiability.  All terms in the ODE have positive sign---no growing
exponentials appear anywhere in the computation, guaranteeing
unconditional numerical stability for arbitrary optical depth.  We
derive the complete system of ODEs (reflection~$\mat{R}$,
transmission~$\mat{T}$, beam source~$\vect{s}$) for both forward and
backward sweeps, present the boundary-condition solve, and discuss the
Kvaerno5 integration via JAX and diffrax.  This report extends
Sections~3--4 of the Pythonic-DISORT Comprehensive
Documentation~\cite{HP2024doc} and assumes familiarity with that
document.
\end{abstract}

\tableofcontents
\vspace{1em}

%======================================================================
\section{Problem Statement}\label{sec:problem}
%======================================================================

We solve the one-dimensional radiative transfer equation (RTE) for the
diffuse intensity $u(\tau, \mu, \varphi)$ in a plane-parallel atmosphere
of total optical depth $\tbot > 0$.  The mathematical setting is
identical to Sections~1--2 of the Comprehensive
Documentation~\cite{HP2024doc}; we restate the key equations here for
self-containment.

%----------------------------------------------------------------------
\subsection{The Radiative Transfer Equation}
%----------------------------------------------------------------------

The RTE for diffuse intensity is (cf.\ \cite[eq.~(1)]{HP2024doc})
\begin{equation}\label{eq:rte}
  \mu \frac{\partial u}{\partial \tau}
  = u(\tau, \mu, \varphi)
  - \frac{\omega(\tau)}{4\pi}\int_{-1}^{1}\!\int_0^{2\pi}
    p(\mu, \varphi; \mu', \varphi'; \tau)\,
    u(\tau, \mu', \varphi')\dd\varphi'\dd\mu'
  - Q(\tau, \mu, \varphi),
\end{equation}
where $\tau \in [0, \tbot]$ is optical depth (increasing downward),
$\mu \in [-1,1]$ is the cosine of the polar angle ($\mu > 0$ upward,
$\mu < 0$ downward), $\omega(\tau) \in [0,1)$ is the single-scattering
albedo, $p$ is the scattering phase function, and
\begin{equation}\label{eq:beam-source}
  Q(\tau, \mu, \varphi)
  = \frac{\omega(\tau)\, I_0}{4\pi}\,
    p(\mu, \varphi; -\mu_0, \varphi_0; \tau)\,
    e^{-\tau/\mu_0}
\end{equation}
is the single-scattering source from the collimated solar beam of
intensity $I_0$ incident at angles $(\mu_0, \varphi_0)$ with
$\mu_0 \in (0,1]$.

%----------------------------------------------------------------------
\subsection{Fourier Decomposition and Quadrature Discretisation}
%----------------------------------------------------------------------

Following Section~3.2 of~\cite{HP2024doc}, we expand in azimuthal
Fourier modes
\begin{equation}\label{eq:fourier}
  u(\tau, \mu, \varphi)
  \approx \sum_{m=0}^{M-1} u^m(\tau, \mu)\,\cos\bigl(m(\varphi_0 - \varphi)\bigr),
\end{equation}
where $M = \texttt{NFourier}$.  Each mode
satisfies an independent RTE (cf.\ \cite[Section~3.4]{HP2024doc}):
\begin{equation}\label{eq:rte-mode}
  \mu\,\frac{\dd u^m}{\dd\tau}
  = u^m(\tau, \mu)
  - \omega(\tau)\!\int_{-1}^{1} D^m(\tau, \mu, \mu')\, u^m(\tau, \mu')\dd\mu'
  - Q^m(\tau, \mu),
\end{equation}
where $D^m$ is the azimuthal-mode phase-function kernel
\begin{equation}\label{eq:Dm}
  D^m(\tau, \mu, \mu')
  = \frac{1}{2}\sum_{\ell=m}^{L-1}
    (2\ell+1)\,\frac{(\ell-m)!}{(\ell+m)!}\,
    g_\ell(\tau)\, P_\ell^m(\mu)\, P_\ell^m(\mu'),
\end{equation}
with $g_\ell(\tau)$ the Legendre moments of the phase function and
$P_\ell^m$ the associated Legendre polynomials.  The beam source for
mode~$m$ is
\begin{equation}\label{eq:Qm}
  Q^m(\tau, \mu) = X^m(\tau, \mu)\, e^{-\tau/\mu_0},
  \qquad
  X^m(\tau, \mu) = \frac{\omega(\tau)\, I_0\,(2 - \delta_{0m})}{4\pi}
  \sum_{\ell=m}^{L-1}
    (2\ell+1)\,\frac{(\ell-m)!}{(\ell+m)!}\,
    g_\ell(\tau)\, P_\ell^m(-\mu_0)\, P_\ell^m(\mu).
\end{equation}

Applying double Gauss--Legendre quadrature with $N$ points per hemisphere
(Section~3.1 of~\cite{HP2024doc}) discretises \eqref{eq:rte-mode} into
the $2N$-dimensional ODE system.  We split the intensity into upward and
downward half-vectors:
\[
  \vect{u}^+_i = u^m(\tau, \mu_i), \quad
  \vect{u}^-_i = u^m(\tau, -\mu_i), \quad
  i = 1,\ldots,N,
\]
where $\mu_i > 0$ are the quadrature abscissae and $w_i$ the
corresponding weights.  In matrix form (cf.\ \cite[Section~3.4]{HP2024doc}):
\begin{equation}\label{eq:2N-system}
  \frac{\dd}{\dd\tau}
  \begin{pmatrix} \vect{u}^+ \\ \vect{u}^- \end{pmatrix}
  = \underbrace{
    \begin{pmatrix} -\mat{\alpha}(\tau) & -\mat{\beta}(\tau) \\
                     \mat{\beta}(\tau)  &  \mat{\alpha}(\tau) \end{pmatrix}
  }_{\displaystyle \mat{A}(\tau)}
  \begin{pmatrix} \vect{u}^+ \\ \vect{u}^- \end{pmatrix}
  +
  \begin{pmatrix} -\tilde{\vect{Q}}^+(\tau) \\
                   \tilde{\vect{Q}}^-(\tau) \end{pmatrix},
\end{equation}
where the $N\!\times\!N$ coefficient blocks are
\begin{equation}\label{eq:alpha-beta}
  \mat{\alpha}(\tau)
  = \mat{M}^{-1}\!\bigl(\omega(\tau)\,\mat{D}^+\!\mat{W} - \mat{I}\bigr),
  \qquad
  \mat{\beta}(\tau)
  = \mat{M}^{-1}\omega(\tau)\,\mat{D}^-\!\mat{W},
\end{equation}
with diagonal matrices $\mat{M} = \diag(\mu_1,\ldots,\mu_N)$,
$\mat{W} = \diag(w_1,\ldots,w_N)$, and
\begin{equation}\label{eq:Dpm}
  D^+_{ij} = D^m(\tau, \mu_i, \mu_j),
  \qquad
  D^-_{ij} = D^m(\tau, \mu_i, -\mu_j).
\end{equation}
The scaled beam-source vectors are
\begin{equation}\label{eq:Qtilde}
  \tilde{Q}^\pm_i(\tau) = \frac{1}{\mu_i}\, Q^m(\tau, \pm\mu_i).
\end{equation}

\begin{remark}[Implementation: pre-computed Legendre tensors]
In the implementation, the user supplies a callable
$\texttt{Leg\_coeffs\_func}(\tau) \to (g_0(\tau), g_1(\tau), \ldots,
g_{L-1}(\tau))$ returning the Legendre moments of the phase function at
each $\tau$.  The associated Legendre polynomials $P_\ell^m(\mu_i)$ at
the quadrature points are pre-computed once at setup time using
\texttt{scipy.special} (not JAX-traceable) and stored as JAX arrays.
The matrices $\mat{D}^\pm$ are then assembled inside the ODE vector
field via \texttt{jnp.einsum} contractions over these pre-computed
tensors and the $\tau$-dependent coefficients $g_\ell(\tau)$, making
the entire integration fully JAX-traceable and autodifferentiable.
The factor $\omega(\tau)$ is applied separately, allowing $\omega(\tau)$
and the phase function to vary independently.
\end{remark}

%----------------------------------------------------------------------
\subsection{Differences from PythonicDISORT}\label{sec:differences}
%----------------------------------------------------------------------

The following features of PythonicDISORT~\cite{HP2024doc} are
\emph{not} implemented in the current Riccati solver.

\begin{enumerate}[label=(\roman*)]
  \item \textbf{Delta-$M$ scaling} (Section~3.3 of~\cite{HP2024doc}):
    The truncation of the forward-scattering peak is not applied.
    The full phase function is used directly.

  \item \textbf{Nakajima--Tanaka (NT) corrections}: Not applied.

  \item \textbf{Eigendecomposition}: PythonicDISORT diagonalises the
    constant coefficient matrix $\mat{A}$ in each homogeneous layer
    (Section~3.5) and constructs the general solution analytically
    (Section~3.6).  The Riccati solver replaces this entirely with
    ODE integration, requiring no eigendecomposition.

  \item \textbf{Stamnes--Conklin substitution} (Section~3.7.2): This
    algebraic trick eliminates growing exponentials from the boundary
    system.  The Riccati formulation achieves the same goal
    structurally---all ODE terms have positive sign---so no substitution
    is needed.

  \item \textbf{Isotropic internal source}: Only the collimated beam
    source is handled; the isotropic source $s(\tau)$ (polynomial in
    $\tau$, Section~3.6.2 of~\cite{HP2024doc}) is not yet supported.

  \item \textbf{Multi-layer piecewise-constant model}: PythonicDISORT
    assumes $\omega$ and $g_\ell$ are constant within each layer
    (Section~4).  The Riccati solver treats them as continuous functions
    of $\tau$, eliminating discretisation error from the piecewise
    approximation.

  \item \textbf{Non-ToA depth evaluation}: Currently only the intensity
    at $\tau = 0$ (top of atmosphere) is returned.
\end{enumerate}

The boundary-condition formulation (\cref{sec:bc}) and the quadrature /
Fourier setup are shared with PythonicDISORT.

%======================================================================
\section{Invariant-Imbedding Riccati Formulation}\label{sec:riccati}
%======================================================================

Rather than solving the $2N$-dimensional initial/boundary-value
problem~\eqref{eq:2N-system} with eigendecomposition, we reformulate it
as an initial-value problem for $N\!\times\!N$ \emph{reflection} and
\emph{transmission} operators using the invariant-imbedding
principle~\cite{ambarzumian1943,bellman1960}.

The idea is to build the slab incrementally: starting from a slab of
zero thickness, we add infinitesimal layers and track how the
slab's bulk reflection and transmission operators evolve.  This yields
a matrix Riccati ODE for reflection $\mat{R}(\sigma)$ and a linear
companion ODE for transmission $\mat{T}(\sigma)$, where $\sigma$ is
the accumulated slab thickness.

%----------------------------------------------------------------------
\subsection{Scattering-Operator Representation}\label{sec:scat-ops}
%----------------------------------------------------------------------

Consider a slab occupying $[\tau_a, \tau_b]$ with $\tau_a < \tau_b$.
Its input--output relation for diffuse intensity is described by four
$N\!\times\!N$ operators:
\begin{align}
  \vect{I}^+(\tau_a) &= \mat{R}_{\mathrm{up}}\,\vect{I}^-(\tau_a)
    + \mat{T}_{\mathrm{up}}\,\vect{I}^+(\tau_b) + \vect{s}_{\mathrm{up}},
    \label{eq:scat-up} \\
  \vect{I}^-(\tau_b) &= \mat{T}_{\mathrm{down}}\,\vect{I}^-(\tau_a)
    + \mat{R}_{\mathrm{down}}\,\vect{I}^+(\tau_b) + \vect{s}_{\mathrm{down}},
    \label{eq:scat-down}
\end{align}
where:
\begin{itemize}[nosep]
  \item $\mat{R}_{\mathrm{up}}$: reflection operator for downward
    illumination from above (returns upward intensity at the top),
  \item $\mat{T}_{\mathrm{up}}$: transmission operator for upward
    illumination from below (transmits to the top),
  \item $\mat{R}_{\mathrm{down}}$: reflection operator for upward
    illumination from below (returns downward intensity at the bottom),
  \item $\mat{T}_{\mathrm{down}}$: transmission operator for downward
    illumination from above (transmits to the bottom),
  \item $\vect{s}_{\mathrm{up}}$, $\vect{s}_{\mathrm{down}}$: beam-source
    contributions (upward at top, downward at bottom) generated within
    the slab by the collimated beam.
\end{itemize}

%----------------------------------------------------------------------
\subsection{Forward Sweep: Building from the Bottom}\label{sec:forward}
%----------------------------------------------------------------------

The forward sweep constructs $\mat{R}_{\mathrm{up}}$,
$\mat{T}_{\mathrm{up}}$, and $\vect{s}_{\mathrm{up}}$ for the full slab
$[0, \tbot]$ by ``growing'' the slab upward from the bottom.

\subsubsection{Integration variable}

We introduce the thickness variable $\sigma \in [0, \tbot]$, where
$\sigma$ represents the thickness of the slab built so far (starting
from the bottom).  The physical optical depth corresponding to the
top of the current slab is
\begin{equation}\label{eq:tau-sigma-fwd}
  \tau = \tbot - \sigma.
\end{equation}
At $\sigma = 0$ the slab has zero thickness (no scattering); at
$\sigma = \tbot$ the full atmosphere has been built.

\subsubsection{Infinitesimal layer addition}

Consider a slab of thickness $\sigma$ with known operators
$(\mat{R}, \mat{T}, \vect{s})$.  We add an infinitesimal layer of
thickness $\dd\sigma$ at the top (i.e.\ at optical depth
$\tau = \tbot - \sigma$).  The coefficient matrix of this thin layer
is evaluated at $\tau = \tbot - \sigma$, giving blocks
$\mat{\alpha}(\tbot - \sigma)$ and $\mat{\beta}(\tbot - \sigma)$.

The infinitesimal layer's own reflection and transmission, to first
order in $\dd\sigma$, are:
\begin{alignat}{2}
  r &= \mat{\beta}\dd\sigma + O(\dd\sigma^2),
  &\qquad
  t &= \mat{I} + \mat{\alpha}\dd\sigma + O(\dd\sigma^2).
  \label{eq:inf-layer}
\end{alignat}
Here $\mat{\beta}$ couples upward and downward streams (off-diagonal
blocks of $\mat{A}$), producing reflection, while $\mat{\alpha}$
modifies same-direction propagation (diagonal blocks), producing
transmission.

\subsubsection{Derivation of the Riccati ODE for \texorpdfstring{$\mat{R}$}{R}}

Adding the thin layer on top of the existing slab, the new reflection
operator $\mat{R}(\sigma + \dd\sigma)$ satisfies the composition rule
(Redheffer star product to first order):
\begin{equation}\label{eq:R-composition}
  \mat{R}(\sigma + \dd\sigma)
  = r + t\,\mat{R}\,(\mat{I} - r\,\mat{R})^{-1}\,t
  \approx r + t\,\mat{R}\,t + O(\dd\sigma^2).
\end{equation}
Substituting \eqref{eq:inf-layer}:
\begin{align}
  \mat{R}(\sigma + \dd\sigma)
  &= \mat{\beta}\dd\sigma
     + (\mat{I} + \mat{\alpha}\dd\sigma)\,\mat{R}\,
       (\mat{I} + \mat{\alpha}\dd\sigma) + O(\dd\sigma^2) \nonumber \\
  &= \mat{R} + \bigl(\mat{\alpha}\mat{R}
     + \mat{R}\mat{\alpha} + \mat{\beta}\bigr)\dd\sigma
     + O(\dd\sigma^2).
\end{align}
However, this misses the term $\mat{R}\mat{\beta}\mat{R}$ that arises
from the multiple-reflection series: light reflected upward by the slab
($\mat{R}$) is partially reflected back down by the thin layer ($r =
\mat{\beta}\dd\sigma$), then reflected up again by the slab
($\mat{R}$), contributing $\mat{R}\,\mat{\beta}\,\mat{R}\dd\sigma$ to
first order.  Including this:
\begin{equation}\label{eq:R-increment}
  \mat{R}(\sigma + \dd\sigma)
  = \mat{R} + \bigl(\mat{\alpha}\mat{R}
    + \mat{R}\mat{\alpha}
    + \mat{R}\mat{\beta}\mat{R}
    + \mat{\beta}\bigr)\dd\sigma
    + O(\dd\sigma^2).
\end{equation}
Taking the limit $\dd\sigma \to 0$:
\begin{equation}\label{eq:riccati-R}
  \boxed{
    \frac{\dd\mat{R}}{\dd\sigma}
    = \mat{\alpha}\mat{R} + \mat{R}\mat{\alpha}
      + \mat{R}\mat{\beta}\mat{R} + \mat{\beta},
    \qquad \mat{R}(0) = \mat{0}.
  }
\end{equation}
This is a \emph{matrix Riccati equation}.  The initial condition
$\mat{R}(0) = \mat{0}$ reflects the fact that a slab of zero thickness
has zero reflectance.

\begin{remark}[Positive signs]\label{rem:positive}
All four terms in \eqref{eq:riccati-R} have positive sign.  Since
$\mat{\beta}$ has non-negative entries (it is proportional to
$\omega \mat{D}^-\mat{W}$), and $\mat{R}$ starts at zero and grows,
$\mat{R}(\sigma)$ remains bounded and non-negative.  No term of the
form $e^{+\lambda\tau}$ with $\lambda > 0$ ever appears---this is the
structural guarantee of numerical stability for thick atmospheres.
\end{remark}

\subsubsection{Derivation of the Companion ODE for \texorpdfstring{$\mat{T}$}{T}}

The transmission operator $\mat{T}(\sigma)$ maps upward intensity
entering at the bottom of the slab to upward intensity exiting at the
top.  When the thin layer $\dd\sigma$ is added at the top, light
transmitted through the existing slab ($\mat{T}$) then passes through
the thin layer.  To first order, the thin layer modifies the transmitted
beam by $t = \mat{I} + \mat{\alpha}\dd\sigma$ and allows a single
back-reflection cycle ($\mat{T} \to$ slab reflects $\mat{R} \to$ thin
layer reflects $\mat{\beta}\dd\sigma \to$ transmit up):
\begin{align}
  \mat{T}(\sigma + \dd\sigma)
  &= t\,(\mat{I} + r\,\mat{R})\,\mat{T} + O(\dd\sigma^2) \nonumber \\
  &= (\mat{I} + \mat{\alpha}\dd\sigma)\,
     (\mat{I} + \mat{\beta}\mat{R}\dd\sigma)\,\mat{T}
     + O(\dd\sigma^2) \nonumber \\
  &= \mat{T} + (\mat{\alpha} + \mat{\beta}\mat{R})\,\mat{T}\dd\sigma
     + O(\dd\sigma^2).
\end{align}
However, noting that $\mat{T}$ transmits from the bottom \emph{up
through} the slab, and the thin layer is added at the top, the correct
composition is:
\begin{equation}
  \mat{T}(\sigma + \dd\sigma)
  = t\,\mat{T} + t\,\mat{R}\,r\,\mat{T} + O(\dd\sigma^2)
  = \mat{T} + \mat{T}\,(\mat{\alpha} + \mat{\beta}\mat{R})\dd\sigma
    + O(\dd\sigma^2).
\end{equation}
Here the factor ordering reflects that light first enters the slab
from below (acted on by the old operators) and then exits through the
new layer at the top.  Taking the limit:
\begin{equation}\label{eq:companion-T}
  \boxed{
    \frac{\dd\mat{T}}{\dd\sigma}
    = \mat{T}\,(\mat{\alpha} + \mat{\beta}\mat{R}),
    \qquad \mat{T}(0) = \mat{I}.
  }
\end{equation}
The initial condition $\mat{T}(0) = \mat{I}$ says that a zero-thickness
slab transmits perfectly.  Equation~\eqref{eq:companion-T} is linear
in $\mat{T}$ (given $\mat{R}$), so it is a companion to the nonlinear
Riccati~\eqref{eq:riccati-R}.

\subsubsection{Derivation of the Beam-Source Companion ODE for \texorpdfstring{$\vect{s}$}{s}}

The source vector $\vect{s}_{\mathrm{up}}(\sigma)$ represents the upward
intensity at the top of the slab generated by the collimated beam
scattering within the slab, in the absence of any diffuse illumination
from above or below ($\vect{I}^-(\tau_a) = \vect{0}$,
$\vect{I}^+(\tau_b) = \vect{0}$).

At physical depth $\tau = \tbot - \sigma$, the beam source
vectors~\eqref{eq:Qtilde} contribute:
\begin{align}\label{eq:q-up-down}
  \vect{q}_{\mathrm{up}}(\tau) &= \mat{M}^{-1} X^m(\tau, \mu_i)
    \quad\text{(upward source at depth $\tau$)}, \\
  \vect{q}_{\mathrm{down}}(\tau) &= \mat{M}^{-1} X^m(\tau, -\mu_i)
    \quad\text{(downward source at depth $\tau$)}.
\end{align}

When we add a thin layer $\dd\sigma$ at the top, it generates an
upward source $\vect{q}_{\mathrm{up}}(\tbot - \sigma)\dd\sigma$ directly
and a downward source $\vect{q}_{\mathrm{down}}(\tbot - \sigma)\dd\sigma$
that reflects off the existing slab.  The total increment to
$\vect{s}$ is:
\begin{equation}
  \vect{s}(\sigma + \dd\sigma)
  = \vect{s}
    + \bigl[(\mat{\alpha} + \mat{R}\mat{\beta})\,\vect{s}
      + \mat{R}\,\vect{q}_{\mathrm{down}}(\tbot - \sigma)
      + \vect{q}_{\mathrm{up}}(\tbot - \sigma)\bigr]\dd\sigma
    + O(\dd\sigma^2).
\end{equation}
The $(\mat{\alpha} + \mat{R}\mat{\beta})\vect{s}$ term accounts for the
existing source being modified as it passes through the new layer and
reflects between the layer and slab.  Taking the limit:
\begin{equation}\label{eq:companion-s}
  \boxed{
    \frac{\dd\vect{s}}{\dd\sigma}
    = (\mat{\alpha} + \mat{R}\mat{\beta})\,\vect{s}
      + \mat{R}\,\vect{q}_1(\sigma) + \vect{q}_2(\sigma),
    \qquad \vect{s}(0) = \vect{0},
  }
\end{equation}
where for the \textbf{forward sweep}:
\begin{equation}\label{eq:q-mapping-fwd}
  \vect{q}_1(\sigma) = \vect{q}_{\mathrm{down}}(\tbot - \sigma),
  \qquad
  \vect{q}_2(\sigma) = \vect{q}_{\mathrm{up}}(\tbot - \sigma).
\end{equation}
The mapping in \eqref{eq:q-mapping-fwd} reflects the physical picture:
$\vect{q}_1$ is the downward source at the current top of the slab
(which reflects off the slab via $\mat{R}$ to become upward), and
$\vect{q}_2$ is the upward source generated directly.

%----------------------------------------------------------------------
\subsection{Backward Sweep: Building from the Top}\label{sec:backward}
%----------------------------------------------------------------------

To solve the boundary-condition system (\cref{sec:bc}), we also need the
``downward'' operators $\mat{R}_{\mathrm{down}}$,
$\mat{T}_{\mathrm{down}}$, and $\vect{s}_{\mathrm{down}}$.  These are
obtained by a backward sweep that builds the slab from the top downward.

\subsubsection{Integration variable}

We again use $\sigma \in [0, \tbot]$, but now $\sigma$ represents the
thickness accumulated from the top.  The physical depth at the bottom
of the current slab is
\begin{equation}\label{eq:tau-sigma-bwd}
  \tau = \sigma.
\end{equation}

\subsubsection{ODEs for the backward sweep}

By the same infinitesimal-layer argument, adding a layer at the bottom
(at depth $\tau = \sigma$) and tracking downward reflection/transmission
yields:
\begin{align}
  \frac{\dd\mat{R}_{\mathrm{down}}}{\dd\sigma}
  &= \mat{\alpha}(\sigma)\,\mat{R}_{\mathrm{down}}
     + \mat{R}_{\mathrm{down}}\,\mat{\alpha}(\sigma)
     + \mat{R}_{\mathrm{down}}\,\mat{\beta}(\sigma)\,\mat{R}_{\mathrm{down}}
     + \mat{\beta}(\sigma),
  \label{eq:riccati-R-bwd} \\
  \frac{\dd\mat{T}_{\mathrm{down}}}{\dd\sigma}
  &= \mat{T}_{\mathrm{down}}\,\bigl(\mat{\alpha}(\sigma)
     + \mat{\beta}(\sigma)\,\mat{R}_{\mathrm{down}}\bigr),
  \label{eq:companion-T-bwd} \\
  \frac{\dd\vect{s}_{\mathrm{down}}}{\dd\sigma}
  &= \bigl(\mat{\alpha}(\sigma)
     + \mat{R}_{\mathrm{down}}\,\mat{\beta}(\sigma)\bigr)\,
     \vect{s}_{\mathrm{down}}
     + \mat{R}_{\mathrm{down}}\,\vect{q}_1^{\mathrm{bwd}}(\sigma)
     + \vect{q}_2^{\mathrm{bwd}}(\sigma),
  \label{eq:companion-s-bwd}
\end{align}
with identical initial conditions:
$\mat{R}_{\mathrm{down}}(0) = \mat{0}$,
$\mat{T}_{\mathrm{down}}(0) = \mat{I}$,
$\vect{s}_{\mathrm{down}}(0) = \vect{0}$.

The source mapping for the \textbf{backward sweep} is:
\begin{equation}\label{eq:q-mapping-bwd}
  \vect{q}_1^{\mathrm{bwd}}(\sigma) = \vect{q}_{\mathrm{up}}(\sigma),
  \qquad
  \vect{q}_2^{\mathrm{bwd}}(\sigma) = \vect{q}_{\mathrm{down}}(\sigma).
\end{equation}
This is the mirror of \eqref{eq:q-mapping-fwd}: in the backward sweep,
the upward source at the current bottom reflects off the growing
downward-facing slab (via $\mat{R}_{\mathrm{down}}$), and the downward
source propagates directly.

\begin{remark}[Symmetry for homogeneous slabs]
For a homogeneous atmosphere ($\omega$ and $g_\ell$ independent of
$\tau$), the coefficient functions satisfy $\mat{\alpha}(\tbot - \sigma)
= \mat{\alpha}(\sigma)$ and $\mat{\beta}(\tbot - \sigma) =
\mat{\beta}(\sigma)$.  In this case the forward and backward Riccati
equations are identical, so $\mat{R}_{\mathrm{up}} =
\mat{R}_{\mathrm{down}}$ (test~14d in the test suite verifies this).
\end{remark}

%----------------------------------------------------------------------
\subsection{State Representation for Numerical Integration}\label{sec:state}
%----------------------------------------------------------------------

The three ODEs \eqref{eq:riccati-R}, \eqref{eq:companion-T},
\eqref{eq:companion-s} (or their backward analogues) are integrated
simultaneously as a single coupled system.  The state is represented as
a \emph{PyTree} (a nested dictionary of JAX arrays):
\begin{equation}\label{eq:state-pytree}
  \vect{y} =
  \bigl\{\,
    \texttt{R}: \mat{R} \in \R^{N \times N},\;
    \texttt{T}: \mat{T} \in \R^{N \times N},\;
    \texttt{s}: \vect{s} \in \R^{N}
  \,\bigr\},
\end{equation}
with total dimension $2N^2 + N$.  Unlike a flattened state vector,
the PyTree representation preserves the matrix structure and avoids
reshape operations.  The diffrax solver natively supports PyTree states,
applying its ESDIRK stages and step-size control to each leaf array.

The right-hand side function evaluates $\mat{\alpha}(\sigma)$ and
$\mat{\beta}(\sigma)$ once per call and assembles the three derivatives:
\begin{align}
  \frac{\dd\mat{R}}{\dd\sigma}
  &= \mat{\alpha}\mat{R} + \mat{R}\mat{\alpha}
     + \mat{R}\mat{\beta}\mat{R} + \mat{\beta},
  \label{eq:rhs-R} \\
  \frac{\dd\mat{T}}{\dd\sigma}
  &= \mat{T}\,(\mat{\alpha} + \mat{\beta}\mat{R}),
  \label{eq:rhs-T} \\
  \frac{\dd\vect{s}}{\dd\sigma}
  &= (\mat{\alpha} + \mat{R}\mat{\beta})\vect{s}
     + \mat{R}\vect{q}_1 + \vect{q}_2.
  \label{eq:rhs-s}
\end{align}
The initial condition is $\mat{R}(0) = \mat{0}$, $\mat{T}(0) = \mat{I}$,
$\vect{s}(0) = \vect{0}$.

The cost per right-hand side evaluation is dominated by three
$N\!\times\!N$ matrix multiplications ($O(N^3)$ each).

%======================================================================
\section{Boundary-Condition Solve}\label{sec:bc}
%======================================================================

After the forward and backward sweeps, we have six operators for the
full slab $[0, \tbot]$:
$(\mat{R}_{\mathrm{up}}, \mat{T}_{\mathrm{up}}, \vect{s}_{\mathrm{up}})$
from the forward sweep and
$(\mat{R}_{\mathrm{down}}, \mat{T}_{\mathrm{down}}, \vect{s}_{\mathrm{down}})$
from the backward sweep.

%----------------------------------------------------------------------
\subsection{Surface Boundary Condition}
%----------------------------------------------------------------------

At the bottom boundary ($\tau = \tbot$), the upward intensity is
\begin{equation}\label{eq:surface-bc}
  \vect{I}^+(\tbot)
  = \mat{R}_{\mathrm{surf}}\,\vect{I}^-(\tbot)
    + \vect{b}_{\mathrm{pos}}^{\mathrm{eff}},
\end{equation}
where $\mat{R}_{\mathrm{surf}}$ is the $N\!\times\!N$ surface
reflection matrix incorporating the BDRF (Bidirectional Reflectance
Distribution Function) and $\vect{b}_{\mathrm{pos}}^{\mathrm{eff}}$
includes any prescribed upward surface emission $\vect{b}_{\mathrm{pos}}$
plus the beam-surface reflection term.

For a Lambertian surface with albedo $\rho$, the $m=0$ mode has
\begin{equation}
  (\mat{R}_{\mathrm{surf}})_{ij}
  = 2\,\frac{\rho}{\pi}\,\mu_j w_j,
\end{equation}
and the beam reflection term is
\begin{equation}
  \vect{b}_{\mathrm{beam}} = \frac{4\mu_0 I_0}{4\pi}\,
  \frac{\rho}{\pi}\,e^{-\tbot/\mu_0}\,\vect{1}.
\end{equation}

At the top boundary ($\tau = 0$), the downward diffuse intensity is
prescribed:
\begin{equation}\label{eq:top-bc}
  \vect{I}^-(0) = \vect{b}_{\mathrm{neg}}.
\end{equation}

%----------------------------------------------------------------------
\subsection{The \texorpdfstring{$N\!\times\!N$}{NxN} System}
%----------------------------------------------------------------------

Substituting the scattering relations \eqref{eq:scat-up}--\eqref{eq:scat-down}
evaluated for the full slab:
\begin{align}
  \vect{I}^+(0) &= \mat{R}_{\mathrm{up}}\,\vect{b}_{\mathrm{neg}}
    + \mat{T}_{\mathrm{up}}\,\vect{I}^+(\tbot)
    + \vect{s}_{\mathrm{up}}, \label{eq:Iplus-top} \\
  \vect{I}^-(\tbot) &= \mat{T}_{\mathrm{down}}\,\vect{b}_{\mathrm{neg}}
    + \mat{R}_{\mathrm{down}}\,\vect{I}^+(\tbot)
    + \vect{s}_{\mathrm{down}}.  \label{eq:Iminus-bot}
\end{align}

Substituting \eqref{eq:Iminus-bot} into the surface
BC~\eqref{eq:surface-bc}:
\begin{equation}
  \vect{I}^+(\tbot)
  = \mat{R}_{\mathrm{surf}}\bigl(
      \mat{T}_{\mathrm{down}}\,\vect{b}_{\mathrm{neg}}
      + \mat{R}_{\mathrm{down}}\,\vect{I}^+(\tbot)
      + \vect{s}_{\mathrm{down}}\bigr)
    + \vect{b}_{\mathrm{pos}}^{\mathrm{eff}}.
\end{equation}
Rearranging:
\begin{equation}\label{eq:bc-system}
  \boxed{
    \bigl(\mat{I} - \mat{R}_{\mathrm{surf}}\,\mat{R}_{\mathrm{down}}\bigr)\,
    \vect{I}^+(\tbot)
    = \mat{R}_{\mathrm{surf}}\bigl(
        \mat{T}_{\mathrm{down}}\,\vect{b}_{\mathrm{neg}}
        + \vect{s}_{\mathrm{down}}\bigr)
      + \vect{b}_{\mathrm{pos}}^{\mathrm{eff}}.
  }
\end{equation}
This is an $N\!\times\!N$ linear system for $\vect{I}^+(\tbot)$---half
the size of PythonicDISORT's $2N\!\times\!2N$ boundary
system~\cite[Section~3.7]{HP2024doc}.

%----------------------------------------------------------------------
\subsection{Recovery of ToA Intensity}
%----------------------------------------------------------------------

Once $\vect{I}^+(\tbot)$ is found from \eqref{eq:bc-system}, the
upward intensity at the top of atmosphere is recovered via
\eqref{eq:Iplus-top}:
\begin{equation}\label{eq:recovery}
  \vect{I}^+(0)
  = \mat{R}_{\mathrm{up}}\,\vect{b}_{\mathrm{neg}}
    + \mat{T}_{\mathrm{up}}\,\vect{I}^+(\tbot)
    + \vect{s}_{\mathrm{up}}.
\end{equation}
The full intensity at $\tau = 0$ for all azimuthal angles is then
reconstructed from the Fourier modes:
\begin{equation}
  u(0, \mu, \varphi)
  \approx \sum_{m=0}^{M-1}
    u^m(0, \mu)\,\cos\bigl(m(\varphi_0 - \varphi)\bigr).
\end{equation}

The upward diffuse flux at the top of atmosphere is
\begin{equation}\label{eq:flux}
  F_{\mathrm{up}}^{\mathrm{diffuse}}(0)
  = 2\pi \sum_{i=1}^{N} w_i\,\mu_i\, I^+_i(0),
\end{equation}
where $I^+_i(0)$ are the entries of $\vect{I}^+(0)$ from the $m=0$ mode.

%======================================================================
\section{Numerical Integration: Kvaerno5 via JAX and diffrax}\label{sec:kvaerno5}
%======================================================================

%----------------------------------------------------------------------
\subsection{Choice of Solver}
%----------------------------------------------------------------------

The coupled system \eqref{eq:rhs-R}--\eqref{eq:rhs-s} is a stiff ODE:
the eigenvalues of the coefficient matrix $\mat{A}$ span a wide range
(for $\texttt{NQuad} = 8$, $|\lambda_{\max}/\lambda_{\min}| \approx
50$), and the Riccati nonlinearity amplifies this stiffness.  An
explicit method would require impractically small step sizes.

We use the \textbf{Kvaerno5} solver from diffrax~\cite{kidger2022neural},
a JAX-based differential equation library.  Kvaerno5 is a 5-stage
\emph{ESDIRK} (Explicit Singly Diagonal Implicit Runge--Kutta) method
with the following properties:
\begin{itemize}[nosep]
  \item \textbf{Order 5}: fifth-order global accuracy with embedded
    fourth-order error estimate for adaptive step-size control.
  \item \textbf{L-stability}: the stability function satisfies
    $|R(\infty)| = 0$, meaning infinitely stiff components are damped
    completely in a single step.  This is strictly stronger than
    A-stability and essential for the Riccati equation where the
    effective stiffness varies as $\mat{R}$ grows.
  \item \textbf{ESDIRK structure}: all implicit stages share the same
    diagonal coefficient $\gamma$, so only a single LU factorisation
    per step is needed (reused across stages), unlike fully implicit
    methods (e.g.\ Radau~IIA) that require factorising a larger coupled
    system.
  \item \textbf{Adaptive step-size control}: a \texttt{PIDController}
    adjusts the step size using a proportional-integral-derivative
    strategy on the local error estimate, with user-specified tolerances
    $\texttt{rtol} = \texttt{tol}$ and
    $\texttt{atol} = \texttt{tol} \times 10^{-3}$.
\end{itemize}

The state is represented as a PyTree (\cref{sec:state}), which diffrax
handles natively---no flattening or reshaping of the $N\!\times\!N$
matrices is required.

\subsubsection{Why JAX and diffrax}

The solver is implemented in JAX~\cite{jax2018github} for two reasons
beyond performance.  First, the ODE vector field must be
\emph{JAX-traceable}: all operations from the $\tau$-dependent Legendre
coefficients through the $\mat{\alpha}$, $\mat{\beta}$ assembly to the
Riccati right-hand side must be expressed as JAX primitives so that the
computation graph is available for automatic differentiation.  Second,
diffrax provides built-in adjoint strategies
(\texttt{RecursiveCheckpointAdjoint}, etc.)\ that enable differentiation
through the entire ODE solve, which is required for the gradient-based
retrieval discussed in \cref{sec:retrieval}.

Float64 arithmetic (\texttt{jax.config.update("jax\_enable\_x64", True)})
is mandatory: double precision is needed to resolve transmission
operators $\mat{T}$ that decay to $O(10^{-13})$ in thick atmospheres.

%----------------------------------------------------------------------
\subsection{Stability Considerations}\label{sec:stability}
%----------------------------------------------------------------------

\subsubsection{No positive exponents}

As noted in \cref{rem:positive}, the Riccati formulation guarantees
that all ODE terms have positive sign.  This means $\mat{R}(\sigma)$
grows monotonically from $\mat{0}$ toward its steady-state value (the
algebraic Riccati equation solution) without any exponentially growing
transients.  The transmission operator $\mat{T}$ decays monotonically
from $\mat{I}$ toward zero.

This structural property is the Riccati analogue of the Stamnes--Conklin
substitution in PythonicDISORT: both eliminate growing exponentials,
but the Riccati approach does so at the ODE level rather than as an
algebraic post-hoc fix.

\subsubsection{Minimum stream count}

The Riccati algebraic equation (the steady-state of \eqref{eq:riccati-R}
for constant coefficients) is ill-conditioned when $\texttt{NQuad} = 4$
(two streams per hemisphere), with $\|\mat{R}_{\mathrm{stab}}\| \approx
10$.  For $\texttt{NQuad} \geq 6$, $\|\mat{R}_{\mathrm{stab}}\| \approx
0.35$ and the solver converges reliably.

\subsubsection{Step count behaviour}

The adaptive Kvaerno5 solver achieves approximately \textbf{35 forward
steps} on the $\tau = 30$ cloud benchmark (adiabatic liquid water cloud,
$\omega$ varying from 0.85 to 0.96, asymmetry parameter from 0.865 to
0.820) at $\texttt{tol} = 10^{-3}$, for a total of approximately 65
steps across both sweeps.  Remarkably, the step count is nearly
independent of $\texttt{NQuad}$:
\begin{center}
\begin{tabular}{lc}
\toprule
\texttt{NQuad} & Forward steps \\
\midrule
8  & 35 \\
16 & 39 \\
32 & 42 \\
\bottomrule
\end{tabular}
\end{center}
This is because the step size is controlled by the variation of the
optical properties with $\tau$, not by the stiffness of the discretised
system (which the L-stable Kvaerno5 method handles implicitly).

%======================================================================
\section{Summary of the Forward Solver}\label{sec:summary}
%======================================================================

The \texttt{pydisort\_riccati\_jax} solver replaces per-layer
eigendecomposition with invariant-imbedding Riccati integration:
\begin{enumerate}[nosep]
  \item Forward sweep ($\sigma$: $0 \to \tbot$, building from bottom):
    integrate the coupled $(\mat{R}, \mat{T}, \vect{s})$ system
    \eqref{eq:riccati-R}--\eqref{eq:companion-s} to obtain
    $(\mat{R}_{\mathrm{up}}, \mat{T}_{\mathrm{up}}, \vect{s}_{\mathrm{up}})$.
  \item Backward sweep ($\sigma$: $0 \to \tbot$, building from top):
    integrate \eqref{eq:riccati-R-bwd}--\eqref{eq:companion-s-bwd} to
    obtain $(\mat{R}_{\mathrm{down}}, \mat{T}_{\mathrm{down}},
    \vect{s}_{\mathrm{down}})$.
  \item Solve the $N\!\times\!N$ boundary system \eqref{eq:bc-system}
    for $\vect{I}^+(\tbot)$.
  \item Recover ToA intensity via \eqref{eq:recovery}.
  \item Repeat for each Fourier mode $m = 0, \ldots, M-1$.
  \item Assemble full intensity via Fourier synthesis.
\end{enumerate}

The approach is unconditionally stable (no positive exponents),
adaptive (tolerance-controlled step size via L-stable Kvaerno5),
efficient (step count nearly independent of $\texttt{NQuad}$), and
fully JAX-traceable (enabling automatic differentiation for
gradient-based retrieval; see \cref{sec:retrieval}).

%======================================================================
\section{Toward Retrieval: Discrete Adjoints via JAX}\label{sec:retrieval}
%======================================================================

The ultimate goal of \texttt{pydisort\_riccati\_jax} is to serve as the
forward model inside an iterative retrieval of the effective radius
profile $r_e(\tau)$ from top-of-atmosphere radiance observations.  This
section discusses the automatic-differentiation infrastructure that
motivates the JAX implementation and the open questions that remain.

\emph{This section describes work in progress.  The forward solver
(\cref{sec:riccati}--\cref{sec:summary}) is complete and verified; the
retrieval components described below are not yet implemented.}

%----------------------------------------------------------------------
\subsection{The Retrieval Problem}\label{sec:retrieval-problem}
%----------------------------------------------------------------------

Given multi-spectral observations $\vect{y}_{\mathrm{obs}} \in \R^m$ of
upwelling radiance at the top of atmosphere, we seek the vertical
profile of cloud droplet effective radius $r_e(\tau)$ that minimises a
cost function of the form
\begin{equation}\label{eq:cost}
  J(\vect{x})
  = \frac{1}{2}\bigl\|\vect{y}_{\mathrm{obs}} - F(\vect{x})\bigr\|^2_{\mat{S}_\epsilon^{-1}}
  + \frac{1}{2}\bigl\|\vect{x} - \vect{x}_a\bigr\|^2_{\mat{S}_a^{-1}},
\end{equation}
where $\vect{x}$ is the state vector parameterising $r_e(\tau)$,
$F(\vect{x})$ is the forward model (a call to
\texttt{pydisort\_riccati\_jax} composed with a lookup table
$r_e \to (\omega, g_\ell)$), $\mat{S}_\epsilon$ is the measurement
error covariance, $\mat{S}_a$ is the \emph{a~priori} covariance, and
$\vect{x}_a$ is the \emph{a~priori} state~\cite{rodgers2000}.

The Gauss--Newton / Levenberg--Marquardt iteration requires the
Jacobian $\mat{K} = \partial F / \partial \vect{x}$, known in the
atmospheric retrieval literature as the matrix of
\emph{weighting functions}.  For an ill-conditioned retrieval---where the
singular values of $\mat{K}$ span many orders of magnitude, as expected
for a vertically resolved $r_e$ profile---the accuracy of $\mat{K}$
directly controls convergence of the iteration and the fidelity of the
posterior error characterisation (averaging kernels, degrees of freedom
for signal)~\cite{rodgers2000}.

%----------------------------------------------------------------------
\subsection{Automatic Differentiation: Forward and Reverse Mode}
\label{sec:ad}
%----------------------------------------------------------------------

Automatic differentiation (AD) computes exact derivatives of a
computer program by systematic application of the chain
rule~\cite{griewank2008}.  Two modes are available:

\begin{itemize}[nosep]
  \item \textbf{Forward mode} computes one directional derivative
    (Jacobian--vector product) per pass.  For a function
    $F : \R^p \to \R^m$, the full Jacobian requires $p$ forward passes,
    each costing $O(1)$ times the forward evaluation.
  \item \textbf{Reverse mode} computes one adjoint (vector--Jacobian
    product) per pass.  The full Jacobian requires $m$ reverse passes.
    For a scalar cost function ($m = 1$), a single reverse pass yields
    the complete gradient $\nabla_{\!\vect{x}} J$ at a cost of at most
    $\bar{\omega} \cdot \mathrm{OPS}(F)$, where $\bar{\omega} \leq 5$
    (the ``cheap gradient principle'', Proposition~3.1
    of~\cite{griewank2008}).  This bound is \emph{independent of~$p$}.
\end{itemize}

The cost asymmetry between modes determines the preferred AD strategy,
but the relevant comparison depends on what the optimisation algorithm
requires.  Two cases arise in the retrieval setting
(see~\cite{baydin2018} for a general survey):

\begin{enumerate}[nosep]
  \item \textbf{Scalar gradient $\nabla J$.}  Gradient-descent and
    L-BFGS methods need only $\nabla_{\!\vect{x}} J \in \R^p$.  Since
    the cost function $J$ is scalar ($m = 1$), a single reverse-mode
    pass suffices regardless of~$p$.

  \item \textbf{Full Jacobian $\mat{K} = \partial F / \partial \vect{x}$.}
    Gauss--Newton and Levenberg--Marquardt iterations, as well as the
    Rodgers optimal-estimation diagnostics (averaging kernels, degrees
    of freedom for signal~\cite{rodgers2000}), require the full
    $m \times p$ Jacobian of the forward model $F : \R^p \to \R^m$,
    where $m$ is the number of observations (viewing angles, spectral
    bands) and $p$ is the number of profile parameters.  Computing
    $\mat{K}$ requires either $m$ reverse-mode passes (one per
    observation, via \texttt{jax.jacrev}) or $p$ forward-mode passes
    (one per parameter, via \texttt{jax.jacfwd}).  \textbf{Reverse mode
    is preferred when $p > m$; forward mode when $p < m$.}
\end{enumerate}

\subsubsection{Concrete example: adiabatic cloud retrieval}

Consider the adiabatic cloud benchmark of \cref{sec:stability}
($\tbot = 30$, $\omega$ varying 0.85--0.96) with $K \approx 50$
Kvaerno5 steps across both sweeps.  Suppose $m = 10$ observations
(multi-angle or multi-spectral) and let $C_{\mathrm{fwd}}$ denote the
cost of one forward evaluation of \texttt{pydisort\_riccati\_jax},
which scales linearly with~$K$.  Each forward-mode pass costs
$c_f \, C_{\mathrm{fwd}}$ (typically $c_f \approx 2$--$3$); each
reverse-mode pass with checkpointing costs
$c_r \, C_{\mathrm{fwd}}$ ($c_r \approx 3$--$5$, due to segment
recomputation).  The total cost for the full Jacobian is then:
\begin{center}
\begin{tabular}{lcc}
\toprule
Parameterisation & $p$ & Preferred mode \\
\midrule
NN decoder / EOF basis  & 3--5   & forward ($\approx 10\,C_{\mathrm{fwd}}$
                                      vs.\ $\approx 40\,C_{\mathrm{fwd}}$) \\
B-spline basis          & 8--15  & forward or comparable \\
Level-by-level + Tikhonov & 30--50 & reverse ($\approx 40\,C_{\mathrm{fwd}}$
                                      vs.\ $\approx 100\,C_{\mathrm{fwd}}$) \\
\bottomrule
\end{tabular}
\end{center}
The crossover occurs at roughly $p \approx (c_r / c_f)\, m \approx
15$--$20$ for $m = 10$.  In all cases, $C_{\mathrm{fwd}} \propto K$,
so minimising the integration step count (\cref{sec:stability})
directly reduces the cost of every Jacobian computation regardless of
the AD mode chosen.  Each Gauss--Newton iteration requires one Jacobian
plus one forward evaluation; a typical retrieval converges in 5--15
iterations, giving a total cost of $\sim\!10 \times (p\,c_f + 1)
\times C_{\mathrm{fwd}}$ in forward mode or
$\sim\!10 \times (m\,c_r + 1) \times C_{\mathrm{fwd}}$ in reverse.

\subsubsection{JAX traceability}

JAX~\cite{jax2018github} implements both modes (\texttt{jax.jacfwd} for
forward, \texttt{jax.jacrev} for reverse, \texttt{jax.grad} for scalar
reverse) via source transformation of Python/NumPy code.  A function is
\emph{JAX-traceable} if all its operations are expressed as JAX
primitives, so that JAX can construct the computation graph.  This is
why the Riccati vector field uses \texttt{jnp.einsum} over pre-computed
Legendre tensors rather than calling \texttt{scipy.special} at each
step (see the Remark in \cref{sec:problem}).

%----------------------------------------------------------------------
\subsection{Discrete vs.\ Continuous Adjoints}\label{sec:disc-vs-cont}
%----------------------------------------------------------------------

Differentiating through an ODE solver admits two fundamentally different
strategies, often called \emph{discretize-then-optimize} (DTO, discrete
adjoint) and \emph{optimize-then-discretize} (OTD, continuous
adjoint)~\cite{giles2000}.

\subsubsection{Continuous adjoint (OTD)}

One derives the adjoint ODE analytically from the continuous forward
problem.  For $\dd\vect{y}/\dd t = \vect{f}(\vect{y}, \vect{\theta})$,
the adjoint satisfies
$\dd\vect{a}/\dd t = -\vect{a}^\top (\partial\vect{f}/\partial\vect{y})$,
integrated backward in time from a terminal
condition~\cite{pontryagin1962,chen2018neural}.  The resulting adjoint
ODE is then discretized and solved numerically.  The gradient obtained
is exact for the \emph{continuous} cost functional but only approximate
for the discrete cost $J_h$ that was actually computed.

\subsubsection{Discrete adjoint (DTO)}

One first discretizes the forward problem (choosing a specific numerical
scheme), then differentiates through the resulting discrete computation
graph using reverse-mode AD.  The gradient obtained is the exact
derivative of the discrete cost $J_h$ to machine precision.

\subsubsection{Why the distinction matters for ill-conditioned problems}

The two approaches converge to the same continuous sensitivity in the
limit of infinite resolution.  In practice, they diverge for three
reasons:

\begin{enumerate}[nosep]
  \item \textbf{Trajectory reconstruction error.}  The continuous adjoint
    must reconstruct the forward trajectory by backward integration.
    Zhuang et al.~\cite{zhuang2020} proved that the reconstruction
    error is $O(h^{p+1})$ and demonstrated empirically that it causes
    gradient estimation errors that degrade optimisation convergence.

  \item \textbf{Discretization inconsistency.}  The continuous adjoint
    ODE is discretised independently of the forward scheme.  There is
    no guarantee that the resulting discrete adjoint matches the true
    transpose of the forward solver's linearisation.  Sirkes and
    Tziperman~\cite{sirkes1997} showed that this can produce
    non-physical computational modes that corrupt time-dependent
    sensitivity analysis.

  \item \textbf{Stiff system instability.}  Kim et
    al.~\cite{kim2021stiff} demonstrated that the continuous adjoint
    method fails on benchmark stiff ODE systems, with the backward
    adjoint solver diverging due to Newton non-convergence.  The
    Riccati ODE is stiff (the eigenvalue ratio of $\mat{A}$ grows with
    $\texttt{NQuad}$), making this a practical concern.
\end{enumerate}

For an ill-conditioned retrieval, the Gauss--Newton update involves
$(\mat{K}^\top \mat{S}_\epsilon^{-1}\mat{K} + \mat{S}_a^{-1})^{-1}
\mat{K}^\top \mat{S}_\epsilon^{-1}(\vect{y}_{\mathrm{obs}} - F(\vect{x}))$.
When $\mat{K}^\top\mat{K}$ has eigenvalues spanning many orders of
magnitude, even small errors in $\mat{K}$ from an approximate
(continuous) adjoint can project onto the ill-conditioned directions and
produce uncontrolled update steps.  The discrete adjoint avoids this by
computing $\mat{K}$ exactly---the gradient is a guaranteed descent
direction for the actual computed cost.

Sandu~\cite{sandu2006} proved that discrete adjoints of Runge--Kutta
methods preserve the order of accuracy and stiff stability properties
of the forward scheme, giving a formal guarantee that the backward pass
through an L-stable method like Kvaerno5 remains well-behaved for stiff
problems.

\subsubsection{Why JAX rather than the adjoint RTE}

The traditional approach in the radiative transfer community is to
derive the adjoint of the radiative transfer equation analytically
and solve it alongside the forward
equation~\cite{rozanov2006adjoint,hasekamp2005,spurr2001lidort}.
Codes such as LIDORT~\cite{spurr2001lidort} and
SCIATRAN~\cite{rozanov2006adjoint} have invested years of effort in
hand-coded analytic Jacobians.

For the Riccati formulation, deriving the continuous adjoint would
require linearising the coupled nonlinear system
\eqref{eq:riccati-R}--\eqref{eq:companion-s}---including the
$\mat{R}\mat{\beta}\mat{R}$ quadratic term---with respect to the
$\tau$-dependent optical properties, then deriving and implementing the
adjoint of the linearised system for both the forward and backward
sweeps.  This is analytically tractable but labour-intensive and
error-prone.

The JAX + diffrax approach avoids this entirely: the discrete adjoint
is computed automatically by reverse-mode AD through the Kvaerno5
solver, requiring no additional mathematical derivation.  The resulting
gradient is exact (discrete adjoint), not approximate (continuous
adjoint), and the implementation tracks code changes automatically---if
the forward model is modified, the adjoint updates accordingly.

%----------------------------------------------------------------------
\subsection{diffrax Adjoint Strategies}\label{sec:diffrax-adjoints}
%----------------------------------------------------------------------

diffrax~\cite{kidger2022neural} provides several adjoint strategies
for differentiating through \texttt{diffeqsolve}.  The relevant options
for the Riccati solver are:

\begin{description}[style=nextline, leftmargin=2em]
  \item[\texttt{RecursiveCheckpointAdjoint}~(recommended)]
    Discrete adjoint via reverse-mode AD with Stumm--Walther recursive
    binomial checkpointing~\cite{griewank2008}.  Memory cost is
    $O(M \log K)$ where $M = 2N^2 + N$ is the state dimension and
    $K$ is the number of accepted steps.  Gradients are exact (of the
    discrete computation).  This is the default and recommended strategy
    in diffrax.

  \item[\texttt{BacksolveAdjoint}~(not recommended)]
    Continuous adjoint via backward integration of the adjoint ODE.
    Memory cost is $O(M)$, but gradients are approximate.  The diffrax
    documentation states: ``As the checkpointing of
    \texttt{RecursiveCheckpointAdjoint} also gives low memory usage,
    [it] is essentially always preferred.''

  \item[\texttt{ForwardMode}]
    Forward-mode AD through the solver (\texttt{jax.jvp}).  Useful
    when $p \ll m$ (few parameters, many outputs), but for the
    retrieval setting ($p > m$), reverse mode is more efficient.

  \item[\texttt{DirectAdjoint}]
    Both forward and reverse support, but cost scales with
    \texttt{max\_steps} regardless of actual steps taken.  Not practical
    for adaptive solvers.
\end{description}

For the Riccati solver with $K \approx 35$ forward steps and
$N = 4$--$16$, \texttt{RecursiveCheckpointAdjoint} with
$\lceil\log_2 K\rceil \approx 6$ checkpoints is the natural choice.

%----------------------------------------------------------------------
\subsection{Numerical Robustness of the Backward Pass}
\label{sec:adjoint-robustness}
%----------------------------------------------------------------------

While the forward Riccati integration is provably stable (all terms
positive, $\mat{R}$ bounded; \cref{rem:positive}), the discrete adjoint
obtained by differentiating through the Kvaerno5 solver does not
automatically inherit this property.  Two potential sources of concern
are:

\begin{enumerate}
  \item \textbf{Implicit stage differentiation.}  Each Kvaerno5 stage
    solves a nonlinear system via Newton iteration.  The backward pass
    must differentiate through these Newton steps, which involves the
    Jacobian $\partial\vect{f}/\partial\vect{y}$ of the ODE
    right-hand side.  The Riccati RHS contains the quadratic term
    $\mat{R}\mat{\beta}\mat{R}$, whose Jacobian with respect to
    $\mat{R}$ is $(\mat{I} \otimes \mat{R}\mat{\beta}) +
    (\mat{R}\mat{\beta} \otimes \mat{I})$ (via the Kronecker-product
    formula for matrix differential calculus).  Near the steady state,
    $\mat{R}$ is $O(1)$ and this Jacobian is well-behaved, but the
    condition number should be monitored during the retrieval.

  \item \textbf{Sensitivity accumulation over many steps.}  The discrete
    adjoint propagates sensitivity information backward through all
    $K$ accepted steps.  If the forward problem has a large sensitivity
    to parameters at a particular $\tau$-layer (e.g.\ near cloud base
    where $\omega$ changes rapidly), this sensitivity is amplified
    through the chain rule.  The L-stability of Kvaerno5 provides
    damping of stiff modes in the \emph{forward} direction; Sandu's
    result~\cite{sandu2006} guarantees that the discrete adjoint of an
    L-stable RK method inherits this stiff stability, but empirical
    verification for the specific Riccati system is warranted.
\end{enumerate}

A recommended validation experiment, prior to implementing the full
retrieval, is to compute $\nabla_{\!\vect{x}} J$ via
\texttt{jax.grad} through \texttt{pydisort\_riccati\_jax} for a
synthetic $r_e(\tau)$ profile and verify agreement with finite-difference
gradients.  This would confirm both the correctness of the AD pipeline
and the absence of numerical pathologies in the backward pass.

%----------------------------------------------------------------------
\subsection{Profile Parameterisation}\label{sec:parameterisation}
%----------------------------------------------------------------------

The retrieval requires a finite-dimensional parameterisation of the
continuous profile $r_e(\tau)$.  The choice of parameterisation affects
both the conditioning of the inverse problem and the number of
parameters~$p$, which in turn determines whether reverse-mode or
forward-mode AD is more efficient.

Several approaches have been used in the atmospheric retrieval
literature:

\begin{description}[style=nextline, leftmargin=2em]
  \item[Level-by-level with smoothness regularisation]
    The profile is retrieved at each of $p$ discrete $\tau$-levels,
    with a Tikhonov penalty on the first or second derivative to
    suppress oscillatory solutions~\cite{steck2002,dubovik2000}.
    This gives the most general parameterisation but produces the
    largest state vector and the most ill-conditioned $\mat{K}$.

  \item[Basis-function expansion]
    The profile is expanded in a set of $p$ smooth basis functions
    (e.g.\ B-splines, wavelets, or piecewise-linear hat functions on
    a coarse grid~\cite{bowman2006}).  The state vector consists of
    the $p$ expansion coefficients, typically $p = 5$--$20$, which
    dramatically reduces the dimension relative to the full
    level-by-level approach while enforcing smoothness by construction.

  \item[EOF / PCA basis from climatological data]
    Principal components of $r_e$ profiles from a climatological
    ensemble define a data-driven reduced basis~\cite{masiello2012}.
    If the first $k$ EOFs capture $>95\%$ of profile variance, the
    state vector is $k$-dimensional (typically $k = 3$--$10$), with
    the regularisation encoded implicitly by the truncation.

  \item[Neural-network profile generator]
    A neural network maps a low-dimensional latent vector
    ($p = 2$--$5$) to a full profile, trained on a climatology of
    physically consistent profiles.  This differentiable decoder is
    used \emph{inside} the physics-based forward model during the
    retrieval~\cite{gebhard2024}.  Advantages include very low $p$
    (easing ill-conditioning) and the ability to encode physical
    constraints (e.g.\ monotonicity, positivity) in the network
    architecture.

  \item[Optimal estimation with \emph{a~priori} covariance]
    The standard Rodgers~\cite{rodgers2000} framework uses a Gaussian
    prior $\mathcal{N}(\vect{x}_a, \mat{S}_a)$ on the profile,
    retrieved at all levels.  The \emph{a~priori} covariance
    $\mat{S}_a$ acts as a smoothness constraint.  This is
    mathematically equivalent to a Gaussian process prior with a
    specific kernel and is the most widely used approach in operational
    retrievals.
\end{description}

The choice of parameterisation has direct implications for the AD mode,
as discussed in \cref{sec:ad}: for a Gauss--Newton retrieval with
$m \approx 10$ observations, forward-mode AD is cheaper when
$p \lesssim 15$ (small basis, EOF, or NN parameterisation), while
reverse-mode AD becomes advantageous for $p \gtrsim 20$ (level-by-level
retrieval).  If the retrieval uses L-BFGS instead of Gauss--Newton,
only the scalar gradient $\nabla J$ is needed and a single reverse-mode
pass suffices regardless of $p$---at the cost of losing the averaging
kernel and degrees-of-freedom diagnostics that require the full
Jacobian~$\mat{K}$.

%----------------------------------------------------------------------
\subsection{Remaining Steps}\label{sec:remaining}
%----------------------------------------------------------------------

The following components are needed to complete the differentiable
retrieval pipeline:

\begin{enumerate}
  \item \textbf{Differentiable lookup table.}  The mapping
    $r_e(\tau) \to \bigl(\omega(\tau),\, g_\ell(\tau)\bigr)$ via Mie
    theory or a pre-computed table must be JAX-traceable.
    Differentiable interpolation (e.g.\ \texttt{jnp.interp} or a
    spline in JAX) would suffice for a pre-computed table; direct Mie
    computation in JAX is also possible but more involved.

  \item \textbf{End-to-end gradient validation.}  Compute
    $\nabla_{\!\vect{x}} J$ via \texttt{jax.grad} through the full
    forward model (lookup $\to$ Riccati ODE $\to$ BC solve $\to$
    Fourier synthesis $\to$ cost) and verify against finite-difference
    gradients for a synthetic $r_e(\tau)$ profile.

  \item \textbf{Adjoint strategy selection.}  Benchmark
    \texttt{RecursiveCheckpointAdjoint} against \texttt{ForwardMode}
    for the expected parameter count to determine the optimal AD mode.

  \item \textbf{Retrieval experiments.}  Implement the Gauss--Newton /
    Levenberg--Marquardt iteration with the AD-computed Jacobian,
    test convergence on synthetic retrievals with known $r_e(\tau)$
    profiles, and characterise the posterior error via averaging
    kernels.
\end{enumerate}

%======================================================================
% References
%======================================================================
\begin{thebibliography}{99}

% ── Original references ──────────────────────────────────────────────
\bibitem{HP2024doc}
  H.~P.~Le and D.~Hogan,
  \emph{Pythonic-DISORT: Comprehensive Documentation},
  2024.
  \url{https://pythonic-disort.readthedocs.io}

\bibitem{ambarzumian1943}
  V.~A.~Ambartsumian,
  ``On the problem of diffuse reflection of light,''
  \emph{Doklady Akademii Nauk SSSR}, vol.~38, pp.~229--232, 1943.

\bibitem{bellman1960}
  R.~Bellman and R.~Kalaba,
  ``Invariant imbedding, wave propagation, and the WKB approximation,''
  \emph{Proc.\ Nat.\ Acad.\ Sci.}, vol.~46, pp.~1646--1649, 1960.

\bibitem{hairer1996}
  E.~Hairer and G.~Wanner,
  \emph{Solving Ordinary Differential Equations II: Stiff and
  Differential-Algebraic Problems}, 2nd ed.\
  Springer, 1996.

\bibitem{stamnes1984}
  K.~Stamnes and R.~A.~Swanson,
  ``A new look at the discrete ordinate method for radiative
  transfer calculations in anisotropically scattering atmospheres,''
  \emph{J.\ Atmos.\ Sci.}, vol.~45, pp.~387--399, 1988.

% ── JAX and diffrax ──────────────────────────────────────────────────
\bibitem{jax2018github}
  J.~Bradbury, R.~Frostig, P.~Hawkins, M.~J.~Johnson, C.~Leary,
  D.~Maclaurin, G.~Necula, A.~Paszke, J.~VanderPlas, S.~Wanderman-Milne,
  and Q.~Zhang,
  ``JAX: composable transformations of Python+NumPy programs,''
  2018.
  \url{http://github.com/jax-ml/jax}

\bibitem{kidger2022neural}
  P.~Kidger,
  ``On neural differential equations,''
  Ph.D.\ thesis, University of Oxford, 2022.
  arXiv:2202.02435.

% ── Automatic differentiation ────────────────────────────────────────
\bibitem{griewank2008}
  A.~Griewank and A.~Walther,
  \emph{Evaluating Derivatives: Principles and Techniques of
  Algorithmic Differentiation}, 2nd ed.\
  SIAM, 2008.

\bibitem{baydin2018}
  A.~G.~Baydin, B.~A.~Pearlmutter, A.~A.~Radul, and J.~M.~Siskind,
  ``Automatic differentiation in machine learning: a survey,''
  \emph{J.\ Mach.\ Learn.\ Res.}, vol.~18, no.~153, pp.~1--43, 2018.

% ── Discrete vs.\ continuous adjoints ────────────────────────────────
\bibitem{giles2000}
  M.~B.~Giles and N.~A.~Pierce,
  ``An introduction to the adjoint approach to design,''
  \emph{Flow Turbulence Combust.}, vol.~65, pp.~393--415, 2000.

\bibitem{sirkes1997}
  Z.~Sirkes and E.~Tziperman,
  ``Finite difference of adjoint or adjoint of finite difference?''
  \emph{Mon.\ Wea.\ Rev.}, vol.~125, no.~12, pp.~3373--3378, 1997.

\bibitem{sandu2006}
  A.~Sandu,
  ``On the properties of Runge--Kutta discrete adjoints,''
  in \emph{Intl.\ Conf.\ Computational Science (ICCS)},
  LNCS~3994, Springer, 2006, pp.~550--557.

\bibitem{chen2018neural}
  R.~T.~Q.~Chen, Y.~Rubanova, J.~Bettencourt, and D.~K.~Duvenaud,
  ``Neural ordinary differential equations,''
  in \emph{Advances in Neural Information Processing Systems~31
  (NeurIPS)}, 2018.

\bibitem{pontryagin1962}
  L.~S.~Pontryagin, V.~G.~Boltyanskii, R.~V.~Gamkrelidze, and
  E.~F.~Mishchenko,
  \emph{The Mathematical Theory of Optimal Processes},
  Wiley, 1962.

\bibitem{zhuang2020}
  J.~Zhuang, N.~Dvornek, X.~Li, S.~Tatikonda, J.~Papademetris,
  and J.~Duncan,
  ``Adaptive checkpoint adjoint method for gradient estimation in
  neural ODE,''
  in \emph{Intl.\ Conf.\ Machine Learning (ICML)}, 2020.

\bibitem{kim2021stiff}
  S.~Kim, W.~Ji, S.~Deng, Y.~Ma, and C.~Rackauckas,
  ``Stiff neural ordinary differential equations,''
  \emph{Chaos}, vol.~31, 093122, 2021.

% ── Atmospheric retrieval ────────────────────────────────────────────
\bibitem{rodgers2000}
  C.~D.~Rodgers,
  \emph{Inverse Methods for Atmospheric Sounding: Theory and Practice},
  World Scientific, 2000.

\bibitem{rozanov2006adjoint}
  V.~V.~Rozanov,
  ``Adjoint radiative transfer equation and inverse problems,''
  in \emph{Light Scattering Reviews}, Springer Praxis Books,
  pp.~339--392, 2006.

\bibitem{hasekamp2005}
  O.~P.~Hasekamp and J.~Landgraf,
  ``Linearization of vector radiative transfer with respect to aerosol
  properties and its use in satellite remote sensing,''
  \emph{J.\ Geophys.\ Res.}, vol.~110, D04203, 2005.

\bibitem{spurr2001lidort}
  R.~J.~D.~Spurr,
  ``LIDORT and VLIDORT: Linearized pseudo-spherical scalar and vector
  discrete ordinate radiative transfer models for use in remote sensing
  retrieval problems,''
  in \emph{Light Scattering Reviews~3}, Springer, 2008, pp.~229--275.

% ── Profile parameterisation ─────────────────────────────────────────
\bibitem{steck2002}
  T.~Steck,
  ``Methods for determining regularization for atmospheric retrieval
  problems,''
  \emph{Appl.\ Opt.}, vol.~41, no.~9, pp.~1788--1797, 2002.

\bibitem{dubovik2000}
  O.~Dubovik and M.~D.~King,
  ``A flexible inversion algorithm for retrieval of aerosol optical
  properties from Sun and sky radiance measurements,''
  \emph{J.\ Geophys.\ Res.}, vol.~105, pp.~20\,673--20\,696, 2000.

\bibitem{bowman2006}
  K.~W.~Bowman, C.~D.~Rodgers, S.~S.~Kulawik, J.~Worden,
  E.~Sarkissian, L.~Osterman, T.~Steck, M.~Luo, A.~Eldering,
  M.~W.~Shephard, H.~Worden, M.~Lampel, S.~Clough, P.~Brown,
  C.~Rinsland, M.~Gunson, and R.~Beer,
  ``Tropospheric emission spectrometer: retrieval method and error
  analysis,''
  \emph{IEEE Trans.\ Geosci.\ Remote Sens.}, vol.~44, no.~5,
  pp.~1297--1307, 2006.

\bibitem{masiello2012}
  G.~Masiello and C.~Serio,
  ``Inversion for atmospheric thermodynamical parameters of IASI data
  in the principal components space,''
  \emph{Q.~J.~R.\ Meteorol.\ Soc.}, vol.~138, pp.~103--117, 2012.

\bibitem{gebhard2024}
  T.~D.~Gebhard, D.~Angerhausen, B.~S.~Konrad, E.~Alei,
  S.~P.~Quanz, and B.~Sch\"olkopf,
  ``Parameterizing pressure--temperature profiles of exoplanet
  atmospheres with neural networks,''
  \emph{Astron.\ Astrophys.}, vol.~681, A3, 2024.

\end{thebibliography}

\end{document}
